\subsection{The Operating System Mediation Layer}

\texttt{print()}, \texttt{input()}, and \texttt{open()} look like pure Python calls, but every external I/O path eventually becomes a kernel-mediated request from a normal user process. CPython uses runtime wrappers; the operating system owns devices, descriptors, page cache, and permission checks.

\subsubsection{Why Python Does Not Read Disks Directly: System Calls, File Handles, and Kernel Buffers}

Disk blocks and directory metadata are kernel-managed. Python stream methods eventually trigger \texttt{read}/\texttt{write}/\texttt{open}/\texttt{close} system calls (or their platform equivalents) after passing through interpreter buffering layers.

\begin{verbatim}
import sys

msg = "These are not the droids you are looking for."
print(type(msg), type(sys.stdout))
for name in ("write", "flush", "fileno", "buffer"):
    print(name, hasattr(sys.stdout, name))
sys.stdout.write("Serenity now!\n")
sys.stdout.flush()
\end{verbatim}

\subsubsection{File Descriptors vs. Python File Objects: Native Resource Handles Wrapped in Managed Runtime Objects}

POSIX descriptors \texttt{0}/\texttt{1}/\texttt{2} are process-local integer handles, not data. \texttt{sys.stdout} is a full Python object with encoding and buffering semantics. \texttt{fileno()} bridges the models.

\begin{codebox}[]{python}{File Descriptors vs. Python File Objects: Native Resource Handles Wrapped in...}{the01}
import os
import sys

stdout_obj = sys.stdout
stdout_fd = stdout_obj.fileno()
print(stdout_obj, stdout_fd, type(stdout_fd))

text = "Chuck Norris counted to infinity. Twice.\n"
data = text.encode("utf-8")
os.write(stdout_fd, data)
sys.stdout.flush()
\end{codebox}


Low-level \texttt{os.write} accepts \texttt{bytes}; text must be encoded explicitly.

\subsubsection{The Layered \texttt{io} Module Architecture: Understanding Text, Buffered, and Raw Subsystems}

CPython constructs file streams as a vertical stack, not a single syscall wrapper:

\begin{itemize}
    \item \texttt{TextIOWrapper} --- top layer; runs encoding/decoding state machines and universal newline translation; \texttt{read()}/\texttt{write()} traffic in \texttt{str} and typically allocate fresh Unicode objects per operation.
    \item \texttt{BufferedIOBase} implementations (\texttt{BufferedReader}, \texttt{BufferedWriter}, \texttt{BufferedRandom}) --- middle layer; maintain heap-allocated block caches (commonly 8 KiB) to batch user-space requests and reduce kernel entry frequency.
    \item \texttt{FileIO} (\texttt{RawIOBase}) --- bottom layer; wraps the POSIX descriptor (\texttt{int}) and transfers raw bytes with minimal policy.
\end{itemize}

Text mode \texttt{open("path", "r")} builds the full stack. Binary mode omits \texttt{TextIOWrapper} but retains buffering unless disabled. Direct \texttt{os.read(fd, n)} / \texttt{os.write(fd, buf)} bypasses both text translation and block caching---higher throughput potential, but buffer sizing, encoding, and newline semantics become caller responsibilities.

\begin{verbatim}
import sys

print(type(sys.stdout))
print(type(sys.stdout.buffer))
print("Text layer:", sys.stdout)
print("Raw byte layer:", sys.stdout.buffer)
\end{verbatim}

\subsubsection{Buffering Layers: Reducing Expensive Kernel Transitions Through Intermediate Memory Buffers}

User/kernel transitions dominate small I/O costs. Intermediate buffers collect bytes in user space before a syscall. Output may therefore lag behind a \texttt{write()} call until the buffer fills or is flushed.
